\section{Numerical Illustration}
\label{sec:numerical_illustration}

This section presents the numerical illustration of the robust harmonic hosting capacity problem of a phase-balanced three-phase power system. Two case studies are considered on the same test network: first, every harmonic order up to the 50th, including the zero-sequence harmonics, without uncertainty on the current angles, to assess the transformer windings and earthing; and second, the characteristic harmonics up to the 25th with uncertain current angles, to assess the choice of the risk factor and the network data and their uncertainty. All results are presented in per-unit using a 100\,MVA basis.

\pgfplotsset{rhhc/.style={
    axis lines=left, xmajorgrids=true, ymajorgrids=true, grid style={gray!10},
    label style={font=\footnotesize}, tick label style={font=\footnotesize},
    every axis plot/.append style={line width=0.6pt, mark size=1.5pt, black},
    legend style={draw=none, font=\scriptsize, legend columns=-1, at={(0.5,1.02)}, anchor=south},
    legend cell align=left},
    reference/.style={gray!60, line width=1.2pt, mark=none}}

\subsection{Test Network}
\label{sec:test_network}

The test network is the IEEE 30-bus system of~\cite{vanacker_shhc}, with 21 loads as harmonic units under the absolute equality fairness criterion. The per-unit-length parameters of Section~\ref{sec:parameter_uncertainty} are computed with \textsc{LineCableModels.jl}~\cite{lcm_2026} in the exact $\pi$-model~\eqref{eq:pi_series}--\eqref{eq:pi_shunt}. The 33\,kV lines are replaced by 22 underground cables of four sizes, 400 to 800\,mm$^{2}$ aluminum, 18/30(36)\,kV, in touching trefoil at 0.8\,m depth. The twelve 132\,kV lines are aluminum conductors steel reinforced of four sizes on a delta tower under one steel earth wire. Cable screens are solidly bonded at both ends, the earth wire is grounded at every tower, and all circuits are ideally transposed. The transformers carry their vector groups, i.e., YNyn0 for the 132/33\,kV units and YNd11 for the condenser step-up units. The fundamental operating point is that of the nominal parameters. Four network models are compared:
\begin{enumerate}
    \item[(i)] \textit{placeholder}, a single 630\,mm$^{2}$ aluminum cable design for all 33\,kV links, with four uncertain inputs, and the 132\,kV lines of~\cite{vanacker_shhc} without line charging;
    \item[(ii)] \textit{cables}, the four cable sizes with four uncertain inputs, i.e., resistance, insulation thickness, permittivity and earth resistivity, and the same 132\,kV lines;
    \item[(iii)] \textit{cables and lines}, case~(ii) with the 132\,kV lines modeled physically, including their charging; and
    \item[(iv)] \textit{full}, case~(iii) with all 23 physical inputs of Section~\ref{sec:parameter_uncertainty}.
\end{enumerate}

\subsection{Zero-Sequence Harmonics and Transformer Earthing}
\label{sec:zero_sequence_results}

The first case study considers network model~(iii) with every harmonic order from 2 to 50, including the 16 zero-sequence orders, without recombination, as in~\cite{vanacker_hosting_2026}: the unit contributions align in the worst case, so that no angle model of the units and no risk factor are needed. The 50th harmonic, at 2.5\,kHz, lies within the valid range of the cable and line model of Section~\ref{sec:parameter_uncertainty}. Each generator is taken to connect through a YNd step-up with a leakage reactance of 0.12\,pu on a rating of 0.2 times its short-circuit power, at~$X/R = 40$. The synchronous condensers at buses~11 and~13 sit behind their own step-up units, whose delta windings block the zero sequence. % generator zero-sequence data: assumption accepted by Tom Van Acker on 2026-11-11
Five variants differ only in the windings and earthing of the transformers with a 33\,kV winding:
\begin{enumerate*}
    \item[(a)] YNyn0 with solidly earthed 33\,kV stars;
    \item[(b)] YNyn0 with the 33\,kV stars earthed through 19.05\,$\Omega$, which limits an earth fault to 1\,kA, i.e., case~(iii);
    \item[(c)] Dyn11, with the delta on the 132\,kV side and the 33\,kV stars as in~(b);
    \item[(d)] YNd11, with the delta on the 33\,kV side and the 33\,kV star of the condenser step-up unearthed; and
    \item[(e)] YNy0 with every 33\,kV star unearthed.
\end{enumerate*}
The clock numbers keep the same phase shift on every 132/33\,kV path.

\begin{figure}[b]
    \centering
    \begin{tikzpicture}
        \begin{semilogyaxis}[rhhc, scale only axis, width=0.78\columnwidth, height=2.5cm, xmin=1.5, xmax=50.5, ymin=3e-5, ymax=4e-3,
            xtick={2,5,10,15,20,25,30,35,40,45,50}, xlabel={harmonic order $\Harmonic$}, ylabel={$\HostingCapacity_{\Harmonic}$ [pu]},
            every axis plot/.append style={mark size=1.1pt, line width=0.5pt},
            xlabel style={yshift=3pt}, ylabel style={yshift=-3pt},
            legend image code/.code={\draw[mark repeat=2, mark phase=2] plot coordinates {(0cm,0cm) (0.14cm,0cm) (0.28cm,0cm)};},
            legend columns=-1, legend style={at={(0.5,1.02)}, anchor=south, draw=none, inner sep=0pt, font=\scriptsize, /tikz/every even column/.append style={column sep=1.5pt}, /tikz/every odd column/.append style={column sep=-1pt}}]
            \pgfplotsinvokeforeach{3,6,...,48}{\draw[gray!18, line width=3.6pt] (axis cs:#1,3e-5) -- (axis cs:#1,4e-3);}
            \addplot[mark=o, solid]               table[x=h, y=ra] {journals/rhhc/data/zero_sequence.dat};
            \addplot[mark=square*, densely dashed] table[x=h, y=rb] {journals/rhhc/data/zero_sequence.dat};
            \addplot[mark=triangle*, dotted]       table[x=h, y=rc] {journals/rhhc/data/zero_sequence.dat};
            \addplot[mark=diamond*, dashdotted]    table[x=h, y=rd] {journals/rhhc/data/zero_sequence.dat};
            \addplot[mark=x, densely dotted, mark size=1.6pt] table[x=h, y=re] {journals/rhhc/data/zero_sequence.dat};
            \legend{(a)\,14.3, (b)\,13.0, (c)\,13.5, (d)\,12.8, (e)\,12.8}
        \end{semilogyaxis}
    \end{tikzpicture}
    \vspace{-1.5ex}
    \caption{Robust harmonic hosting capacity per harmonic without recombination, earthing variants~(a)--(e), network model~(iii), $d = 4$; legend: total robust hosting capacity~$\sum_{\Harmonic \in \HarmonicSet} \HostingCapacity_{\Harmonic}$ in~$10^{-3}$\,pu; gray: zero-sequence harmonics.}
    \label{fig:zero_sequence}
\end{figure}

Fig.~\ref{fig:zero_sequence} shows that the earthing acts on the zero sequence only. At the positive- and negative-sequence harmonics, the robust hosting capacity of the five variants agrees to within~$10^{-6}$\,pu.

Solid earthing of the 33\,kV stars hosts~42\,\% more zero-sequence current than the earthing through 19.05\,$\Omega$ after robustness. Removing the earthing altogether, variants~(d) and~(e), lowers the robust zero-sequence hosting capacity by only a further 6--8\,\%, as the resistance already decouples the 33\,kV network. A delta on the 132\,kV side, variant~(c), raises it by 15\,\%, as the 132\,kV network no longer shares the zero-sequence path. The effect of the earthing fades with the harmonic order: up to the 24th harmonic, the robust zero-sequence hosting capacity differs by a factor of up to 5.7 between the variants, from the 33rd harmonic on by at most 12\,\%, as the zero-sequence current returns increasingly through the capacitances of the cables and lines rather than the transformer neutrals: per unit injected at a unit bus, the current through the transformer earthing falls from 0.39--0.68 at the 3rd harmonic to at most 0.07 at the 45th. The earthing practice of the distribution network is thus a lever on the zero-sequence hosting capacity of the order of the uncertainty of the network itself.

The uncertain cable parameters cost 9--16\,\% of the zero-sequence hosting capacity, against 6--15\,\% for the positive- and negative-sequence harmonics. The loss concentrates at a few high orders, the 24th and 30th harmonic, where the robust hosting capacity falls to 0.57--0.67 of the nominal one, depending on the earthing. The zero-sequence parameters of the cables and the earthing path together place the zero-sequence resonances, so a deterministic study of the zero-sequence harmonics with nominal cable data misses the harmonics that are most at risk.

\subsection{Harmonics with Uncertain Current Angles}
\label{sec:ieee30_physical}\label{sec:risk_factor_results}

The second case study considers the characteristic harmonics up to the 25th,  with recombination. The current angle of every unit is a random variable, affine in one Beta-distributed random variable per unit, with the angle distributions of~\cite{vanacker_shhc}, so that the angles are independent between the units and fully dependent across the harmonics of a unit. The unit contributions~\eqref{eq:response} recombine at every bus and branch, and the chance constraints~\eqref{eq:cc_compact}, with the exact moments~\eqref{eq:mean_vector}--\eqref{eq:gram_matrix}, limit the result at the risk level~$\varepsilon = 0.05$. The robust hosting capacity is computed for the four network models~(i)--(iv), with the 21 units as diode bridges. First, the choice of the risk factor is assessed on the purely inductive network of~\cite{vanacker_shhc}, without parameter uncertainty, for the diode bridges. 

The Gaussian and Cantelli factors (\S\,\ref{sec:shhc_recap}) are compared on the nominal purely inductive IEEE 30-bus system, i.e., without parameter uncertainty, where the exact~$(1-\varepsilon)$ quantile follows from Monte Carlo simulation. Both act on the exact moments. Fig.~\ref{fig:risk_factor} compares them with the exact quantile for a risk level~$\varepsilon \in \{0.01, 0.02, 0.05, 0.1, 0.2\}$. The Cantelli bound never exceeds the risk level, with violation probabilities of at most 5\,\% of~$\varepsilon$, but costs~2.6--13\,\% of hosting capacity. The Gaussian factor stays within~1.0\,\% of the exact quantile, with a violation probability below the risk level for~$\varepsilon \leq 0.1$ and of~1.07 times the risk level at~$\varepsilon = 0.2$. The Gaussian factor on exact moments is therefore used in the remainder of this section.

\begin{figure}[!b]
    \centering
    \pgfplotsset{rfplot/.style={rhhc, scale only axis, width=0.3\columnwidth, height=2.0cm, xmin=0.008, xmax=0.25,
        xtick={0.01,0.05,0.2}, xticklabels={0.01,0.05,0.2}, xlabel={risk level $\varepsilon$}},
        can/.style={solid, mark=o}, gau/.style={densely dashed, mark=square*}}
    \pgfplotslegendfromname{rflegend}\\[-2pt]
    \subfloat[hosting capacity relative to the exact $(1-\varepsilon)$ quantile]{%
    \begin{tikzpicture}
        \begin{semilogxaxis}[rfplot, ymin=-15, ymax=2, ylabel={difference [\%]},
            legend columns=3, legend to name=rflegend, legend style={font=\scriptsize, /tikz/every even column/.append style={column sep=4pt}}]
            \addplot[reference] coordinates {(0.008,0) (0.25,0)};
            \addplot[can] table[x=eps, y=ddcan] {journals/rhhc/data/risk_factor.dat};
            \addplot[gau] table[x=eps, y=ddgau] {journals/rhhc/data/risk_factor.dat};
            \legend{target, Cantelli, Gaussian}
        \end{semilogxaxis}
    \end{tikzpicture}}\hfill
    \subfloat[violation probability of the binding limit]{%
    \begin{tikzpicture}
        \begin{loglogaxis}[rfplot, ymin=1e-4, ymax=0.3, ylabel={probability [-]}, ytick={1e-3,1e-2,1e-1}, extra y ticks={1.5e-4}, extra y tick labels={0}, extra y tick style={grid=none}]
            \addplot[reference] coordinates {(0.008,0.008) (0.25,0.25)};
            \addplot[can] table[x=eps, y expr={max(\thisrow{vdcan},1.5e-4)}] {journals/rhhc/data/risk_factor.dat};
            \addplot[gau] table[x=eps, y=vdgau] {journals/rhhc/data/risk_factor.dat};
        \end{loglogaxis}
    \end{tikzpicture}}
    \caption{Choice of the risk factor, nominal IEEE 30-bus system, $N = 21$ diode bridges, harmonics up to the 25th: (a) total hosting capacity with the Cantelli and Gaussian factor on exact moments relative to the exact $(1-\varepsilon)$ quantile, and (b) violation probability of the binding individual harmonic voltage limit by Monte Carlo simulation with 20\,000 samples, against the risk level~$\varepsilon$ (gray: target, i.e., the exact quantile in~(a) and a violation probability equal to~$\varepsilon$ in~(b); violation probabilities of zero at the bottom).}
    \label{fig:risk_factor}
\end{figure}

Fig.~\ref{fig:ieee30_physical} gives the results for
\begin{enumerate*}
    \item [(a)] the total harmonic hosting capacity for the four network models, and
    \item [(b)] the per-harmonic robust hosting capacity relative to the nominal stochastic one.
\end{enumerate*}
No case shows a violation on 1024 random points of~$\ParamSet$. At the worst-case points of the final scenario sets of case~(iv), the sampled 95\,\% quantile of the binding limits lies between 0.5\,\% below and 3.0\,\% above the limit, and their violation probability between 0.035 and 0.064, i.e., at most 1.3 times the risk level.

\begin{figure}[!tb]
    \centering
    \pgfplotsset{nmplot/.style={rhhc, scale only axis, width=0.25\columnwidth, height=2.2cm,
        legend style={at={(0.5,1.03)}, anchor=south, draw=none, font=\scriptsize, /tikz/every even column/.append style={column sep=2pt}},
        legend image code/.code={\draw[mark repeat=2, mark phase=2] plot coordinates {(0cm,0cm) (0.15cm,0cm) (0.3cm,0cm)};}},
        mi/.style={solid, mark=o}, mii/.style={densely dashed, mark=square*}, miii/.style={dotted, mark=triangle*}, miv/.style={dashdotted, mark=diamond*},
        res/.style={only marks, mark=|, mark size=2.5pt, line width=0.7pt, forget plot}}
    \subfloat[total]{%
    \begin{tikzpicture}
        \begin{axis}[nmplot, ybar, bar width=4pt, ymin=0, ymax=0.165, ylabel={$\sum_{\Harmonic} \HostingCapacity_{\Harmonic}$ [pu]}, ylabel style={yshift=-4pt},
            symbolic x coords={(i),(ii),(iii),(iv)}, xtick=data, enlarge x limits=0.18, legend columns=2, xlabel={network model},
            legend image code/.code={\draw[#1] (0cm,-0.08cm) rectangle (0.2cm,0.08cm);},
            scaled y ticks=false, yticklabel style={/pgf/number format/fixed, /pgf/number format/precision=2},
            nodes near coords style={font=\tiny, rotate=90, anchor=west, inner sep=1pt}]
            \addplot[black, fill=white] coordinates {((i),0.1330) ((ii),0.1282) ((iii),0.0993) ((iv),0.0993)};
            \addplot[black, fill=gray!60, nodes near coords, point meta=explicit symbolic]
                coordinates {((i),0.1145)[0.86] ((ii),0.1038)[0.81] ((iii),0.0797)[0.80] ((iv),0.0769)[0.78]};
            \legend{nominal, robust}
        \end{axis}
    \end{tikzpicture}}\hfill
    \subfloat[ratio per harmonic]{%
    \begin{tikzpicture}
        \begin{axis}[nmplot, xmin=4, xmax=26, ymin=0.3, ymax=1.05, xtick={5,7,11,13,17,19,23,25}, xticklabels={5,,11,,17,,23,},
            xlabel={harmonic order $\Harmonic$}, ylabel={ratio [-]}, legend columns=4, ylabel style={yshift=-4pt},
            every axis plot/.append style={mark size=1.3pt}]
            \addplot[reference, forget plot] coordinates {(4,1) (26,1)};
            \addplot[mi]   coordinates {(5,0.834) (7,0.896) (11,0.865) (13,0.899) (17,0.861) (19,0.866) (23,0.845) (25,0.854)};
            \addplot[mii]  coordinates {(5,0.906) (7,0.966) (11,0.782) (13,0.851) (17,0.866) (19,0.761) (23,0.747) (25,0.758)};
            \addplot[miii] coordinates {(5,0.884) (7,0.483) (11,0.885) (13,0.623) (17,0.867) (19,0.989) (23,0.929) (25,0.738)};
            \addplot[miv]  coordinates {(5,0.825) (7,0.447) (11,0.874) (13,0.581) (17,0.859) (19,0.937) (23,0.868) (25,0.711)};
            \addplot[res, gray] coordinates {(6.86,0.36) (18.22,0.36) (6.95,0.36) (18.35,0.36)};
            \addplot[res, black] coordinates {(6.27,0.43) (12.32,0.43) (16.86,0.43) (18.21,0.43) (19.12,0.43) (19.65,0.43) (22.72,0.43)};
            \legend{(i), (ii), (iii), (iv)}
        \end{axis}
    \end{tikzpicture}}
    \caption{Nominal stochastic and robust harmonic hosting capacity of the IEEE 30-bus system for the network models~(i)--(iv), harmonics up to the 25th, Gaussian factor on exact moments, $\varepsilon = 0.05$: (a) total, with the robust-to-nominal ratio above the bars, and (b) robust-to-nominal ratio per harmonic, with the parallel resonances at the nominal parameters as ticks at the bottom, gray for~(i) and~(ii), black for~(iii) and~(iv).}
    \label{fig:ieee30_physical}
\end{figure}

With cables only, i.e., case~(ii), the parallel resonances lie at~$\Harmonic = 6.95$ and~18.4. Modeling the 132\,kV lines physically, i.e., case~(iii) moves the first resonance to~$\Harmonic = 6.27$ and adds resonances at~$\Harmonic = 12.32$ and between the 17th and the 23rd harmonic. The robust hosting capacity relative to the nominal one halves at the 7th harmonic, from 0.966 to 0.483, and drops from 0.851 to 0.623 at the 13th harmonic. The network data that enter the robust problem thus matter as much as their uncertainty. A simplified network places the resonances, and with them the worst cases, at other harmonics.

Extending the parameter set from 4 to 23 physical inputs, i.e., case~(iv), lowers the robust hosting capacity by a further~3.5\,\%, mostly at the 5th, 7th, 13th, 19th and 25th harmonic. The uncertainty of the inputs beyond the four dominant ones is therefore not negligible, but it is second order. With~$d = 23$, the dictionary holds 191 sign-predicted corners in place of the~$2^{23}$ vertices. Without them, the method signs off after five iterations at a solution that violates a limit by~0.16\,\% at a sign-predicted corner, whereas 1024 independent random points of~$\ParamSet$ find no violation. With the corners, no violation remains. In high dimension, random sampling, including the a-posteriori check, does not certify robustness, the corners do to first order. The worst corner of the binding limit, the 23rd harmonic at bus~3, is the least damped network, i.e., the coldest cables and lines and the lowest earth resistivity. Moreover, the robust-to-nominal ratio ranges from 0.45 to 0.99 over the harmonics, Fig.~\ref{fig:ieee30_physical}(b). A robust budget is therefore a new allocation, not a uniform derating of the nominal one.

A common simplification scales the fundamental $\pi$-model with the harmonic order, i.e., the resistance with~$\sqrt{\Harmonic}$ and the reactance and susceptance with~$\Harmonic$. Against the physical model, Fig.~\ref{fig:scaling}(a), the rule overestimates the resistance of the lines by a factor of two to three and the reactance of the cables by up to 2.4 at the 25th harmonic, as the screen return lowers their inductance, while the physical cable resistance exceeds the rule by up to 36\,\%. The scaled model thereby moves the parallel resonance near the 23rd harmonic, where the binding limit lies, and halves the driving-point impedance at bus~3 there, Fig.~\ref{fig:scaling}(b). Applied to the nominal 50\,Hz parameters of case~(iv), it lowers the nominal hosting capacity by~21\,\%, yet its allocation over the harmonics violates the limits of the physical model at the nominal parameters. Factor intervals derived from the same box, i.e., the hull of the ratios of the resistance, inductance and capacitance of each family to their nominal values over the box, at the fundamental and at every harmonic frequency, give a robust hosting capacity of 0.0482\,pu, 37\,\% below the physical one, which still violates the limits of the physical model at every one of 1024 random points of~$\ParamSet$. The same factors applied to the nominal frequency response of \textsc{LineCableModels.jl} give 0.0726\,pu, 5.7\,\% below the physical one, with no violation of the physical model at 1024 random points. The frequency dependence of the cable and line parameters is therefore part of the robust problem, and a box of factors on the 50\,Hz parameters does not contain the physical uncertainty. \add{In the zero sequence, the scaling rule fails altogether, as the resistance follows the return path rather than the skin effect of the conductor: up to the 15th harmonic, it rises by only 26\,\% for the cables, whose screens carry the return, and 6.5 times for the lines, whose return runs through the earth, against 3.9 times for~$\sqrt{\Harmonic}$.}

\begin{figure}[!tb]
    \centering
    \pgfplotsset{scplot/.style={rhhc, scale only axis, width=0.29\columnwidth, height=2.2cm, xlabel={harmonic order $\Harmonic$},
        legend columns=2, legend style={at={(0.5,1.03)}, anchor=south, draw=none, font=\scriptsize, /tikz/every even column/.append style={column sep=2pt}},
        legend image code/.code={\draw plot coordinates {(0cm,0cm) (0.3cm,0cm)};}}}
    \subfloat[physical relative to the rule]{%
    \begin{tikzpicture}
        \begin{axis}[scplot, xmin=1, xmax=25, ymin=0.2, ymax=1.5, xtick={1,5,10,15,20,25}, ylabel={ratio [-]}]
            \addplot[reference, forget plot] coordinates {(1,1) (25,1)};
            \addplot[solid, mark=none]          table[x=h, y=Rc] {journals/rhhc/data/scaling_branch.dat};
            \addplot[densely dashed, mark=none] table[x=h, y=Xc] {journals/rhhc/data/scaling_branch.dat};
            \addplot[dotted, line width=0.9pt, mark=none] table[x=h, y=Rl] {journals/rhhc/data/scaling_branch.dat};
            \addplot[dashdotted, mark=none]     table[x=h, y=Xl] {journals/rhhc/data/scaling_branch.dat};
            \legend{$R$ cable, $X$ cable, $R$ line, $X$ line}
        \end{axis}
    \end{tikzpicture}}\hfill
    \subfloat[driving-point impedance at bus~3]{%
    \begin{tikzpicture}
        \begin{semilogyaxis}[scplot, xmin=2, xmax=26, ymin=0.05, ymax=30, xtick={5,10,15,20,25}, ylabel={$|Z_{\Harmonic,33}|$ [pu]}]
            \addplot[gray!60, line width=0.8pt, forget plot] coordinates {(23,0.05) (23,30)};
            \addplot[solid, mark=none]          table[x=h, y=phys]   {journals/rhhc/data/scaling_z.dat};
            \addplot[densely dashed, mark=none] table[x=h, y=scaled] {journals/rhhc/data/scaling_z.dat};
            \legend{physical, scaled}
        \end{semilogyaxis}
    \end{tikzpicture}}
    \caption{Frequency scaling of the fundamental parameters, network model~(iv) at the nominal parameters: (a) series resistance~$R$ and reactance~$X$ of a 500\,mm$^{2}$ cable and a Lion line, physical relative to the scaling rule~$R\sqrt{\Harmonic}$ and~$X\Harmonic$ (gray: rule), and (b) driving-point impedance magnitude at bus~3, physical and scaled (gray: 23rd harmonic, binding limit).}
    \label{fig:scaling}
\end{figure}

\subsection{Computational Tractability}
\label{sec:tractability}

The master problem keeps two variables per harmonic and one cut per scenario, so the effort lies in the oracle. One complete frequency scan of \textsc{LineCableModels.jl} takes about~1.5\,ms per cable design and~2.8--3.6\,ms per overhead-line design, per parameter point and frequency. The full model, with 23 inputs and eight harmonics, requires about 477\,000 such evaluations. With uncertain current angles, every network model converges in two or three iterations to a relative gap below~$5\cdot10^{-4}$, in 34 to 169\,s on 10 threads. The sign-predicted corners reduce the iterations of the full model from five to three and its time by 30\,\%. Without recombination, each earthing variant solves its 49 harmonics in 5.5--7.5\,minutes on 16 threads, in two iterations, to a relative gap below~$10^{-4}$.

\subsection{Comparison with the Literature}
\label{sec:literature_comparison}

With respect to Table~\ref{tab:literature_review}, the robust counterpart solves the 30-bus system in under three minutes with uncertain current angles, and in under eight with 49 harmonics, against under 0.04 minutes for the stochastic problem of~\cite{vanacker_shhc} and 30 minutes for the chance-constrained optimal power flow of~\cite{sun_distribution_2022} on 16 nodes, both with a known network. The studies of the network harmonic impedance~\cite{testa_probabilistic_2002,li_probabilistic_2018,li2020,barakou_parameter_2023} sample or scan the parameters but do not determine a hosting capacity. Whereas an 8\,\% tolerance on the cable capacitance has a minor effect on the harmonic amplification in~\cite{barakou_parameter_2023}, the joint physical inputs cost 14--23\,\% of the hosting capacity here, as the worst cases lie near resonance. The robust methods of~\cite{christakou_voltage_2018,liu_feasible_2023} guarantee feasibility over interval line parameters in linearized models at the fundamental frequency. In this work, the guarantee holds over the exact harmonic network, at the price of an approximate oracle, which the corners and the a-posteriori check control.